\section{Core Components}
\label{sec:components}

This section describes the components whose behaviour the rest of the paper's
claims depend on. We omit components that are conventional --- configuration,
persistence, the storage service --- except where their contract is load-bearing.

\subsection{The registries}

Each family's registry owns typed definitions and their history, and the shapes
differ because the families differ. Prompts hold immutable versions in the
\mlflow{} prompt registry behind aliases; \sys stores the control-plane metadata
about them and does not duplicate the version bodies. Workflows hold editable
drafts promoted to published version rows, with deployment aliases selecting one.
Knowledge bases hold immutable build versions behind an active-version pointer.
Skills hold a mutable current record plus immutable snapshots. Tools hold separate
name-and-version rows with lifecycle status and no live alias. Test sets hold a
version counter plus per-example validity intervals, which is what allows a score
computed last month to be interpreted against the corpus as it stood then.

The uniformity is in the \emph{interface} --- list, read, version, diff --- and not
in the semantics behind it. \tabref{tab:families} is the authority on the latter.

\subsection{The evidence base}

Evidence is the input to every measurement and the reason a release is
reconstructable. Three sources feed it. Test sets are curated example collections
with expectations, versioned so that a score is interpretable later. Production
traces supply real inputs, and trace-to-example capture promotes a production
failure into durable test data --- which is the mechanism that stops the corpus from
drifting away from what production actually does. Assessments carry human and
automated judgements attached to traces.

A design point worth isolating: test sets can be synchronized to \mlflow{}'s native
dataset registry for dataset-level lineage while the relational store remains
authoritative, with a synchronization-state indicator surfaced in the UI. This is a
concrete instance of the sibling-not-proxy relationship --- \sys neither hides
\mlflow{}'s dataset features nor pretends its own copy is a cache.

\subsection{The evaluation framework}

Evaluation composes deterministic scorers with authored LLM judges. Deterministic
scorers are cheap, reproducible, and appropriate for structural properties ---
schema conformance, required-field presence, refusal detection, latency and cost
budgets. Authored judges are created from natural-language instructions through
\mlflow{}'s judge-authoring API, stored as governed artifacts, and referenced by a
\code{Judge.<id>} token anywhere a scorer is accepted: in the refinement gate, in
human review, and in a scorecard.

Because judges are themselves LLM calls, they inherit the biases documented for
LLM-as-judge evaluation~\cite{zheng2023judging,liu2023geval}. \sys's response is
structural rather than corrective: judges are governed artifacts with versions, and
their agreement with human labels is measurable through review queues, so a judge's
trustworthiness is an empirical quantity attached to the judge rather than an
assumption made about it. \appref{app:protocol} specifies how that agreement is to
be measured; \secref{sec:limitations} is explicit that we have not yet measured it.

Human review is built \sys-native over the open-source assessment primitives:
an operator defines a label schema --- pass/fail, categorical, numeric, free text ---
enqueues traces, and reviews them, with answers written back onto each trace as
assessments and expectations. This keeps human labels in the same substrate as
automated ones, which is what makes agreement computable at all.

\subsection{The calibration framework}

%% fig-pipeline.tex -- Figure (single column): inside one refinement job.
%%
%% The internals of the middle stages of Figure fig:chain, at the resolution where
%% the gate semantics become visible. The structural fact the figure exists to
%% make unmissable: the enforced gate governs *candidate advancement*, and the
%% persisted per-version verdict is separate and advisory. They are drawn with
%% different borders so a reader cannot merge them.
%%
%% Rows are placed from the previous row's anchor, so the layout survives a change
%% of wording or of type scale.

\begin{figure}[tbp]
\centering
\begin{cbscaled}[0.82]
\begin{tikzpicture}[x=1cm,y=1cm]

  \def\W{8.06}
  \def\HW{3.78}
  \tikzset{prow/.style  = {anchor=north west, font=\figB, align=left},
           pfull/.style = {prow, text width=\W cm-8pt}}

  % ---- inputs --------------------------------------------------------------
  \node[store, prow, text width=2.30cm] (tr) at (0,0)
    {\textbf{Trace evidence}\\assessments, spans,\\captured examples};
  \node[store, prow, text width=2.30cm] (ts) at ([xshift=2.78cm]tr.north west)
    {\textbf{Test set}\\versioned examples,\\validity intervals};
  \node[store, prow, text width=2.30cm] (bl) at ([xshift=5.56cm]tr.north west)
    {\textbf{Baseline}\\the current \livetarget{}\\and its scores};

  % ---- evidence assembly ---------------------------------------------------
  \node[control, pfull] (ev) at ([yshift=-4.4mm]tr.south west)
    {\textbf{Evidence assembly}\ \ fold traces and examples into one scored
     corpus};
  \foreach \n in {tr,ts,bl} { \draw[flow] (\n.south) -- (\n.south |- ev.north); }

  % ---- diagnosis -----------------------------------------------------------
  \node[control, pfull] (dg) at ([yshift=-3.4mm]ev.south west)
    {\textbf{Diagnosis}\ \ LLM root cause, persisted as a durable artifact};
  \draw[flow] (ev.south) -- (dg.north);

  % ---- optimizer selection -------------------------------------------------
  \node[control, pfull] (op) at ([yshift=-3.4mm]dg.south west)
    {\textbf{Policy selects an optimizer}\ \ by \family and configuration\\[0.4ex]
     \mbox{MetaPrompt}~\dt~\mbox{GEPA}\ \ {\itshape prompt form}\\
     \mbox{SkillMetaPrompt}~\dt~\mbox{BootstrapFewShot}\ \ {\itshape policy}\\
     \mbox{MIPRO}\ \ {\itshape explicit configuration}};
  \draw[flow] (dg.south) -- (op.north);

  % ---- candidate -----------------------------------------------------------
  \node[asset, pfull] (cd) at ([yshift=-3.4mm]op.south west)
    {\textbf{Candidate artifact}\ \ a proposed version, not yet live};
  \draw[flow] (op.south) -- (cd.north);

  % ---- evaluation ----------------------------------------------------------
  \node[control, pfull] (evl) at ([yshift=-3.4mm]cd.south west)
    {\textbf{Evaluation}\ \ scorecard over the assembled corpus\\
     deterministic scorers \dt authored judges as \code{Judge.<id>} \dt
     per-dimension scores against the baseline};
  \draw[flow] (cd.south) -- (evl.north);

  % ---- the two gates, drawn differently -----------------------------------
  \node[govern, prow, text width=\HW cm, line width=1.3pt] (g1)
    at ([yshift=-4.6mm]evl.south west)
    {\textbf{\advgate}\\[0.2ex]
     \textsc{enforced}. Decides whether the job advances to
     \code{candidate\_ready}. A failing gate stops it here.};
  \node[govern, prow, text width=\HW cm, densely dashed] (g2)
    at ([xshift=4.22cm]g1.north west)
    {\textbf{\relverdict}\\[0.2ex]
     \textsc{advisory}. Persisted per version as release evidence. Never blocks
     alias rotation.};
  \draw[flow] ([xshift=0.9cm]g1.north west) -- ++(0,0.46);
  \draw[flow, densely dashed, draw=cbGovern]
    ([xshift=0.9cm]g2.north west) -- ++(0,0.46);

  % ---- human decision + release -------------------------------------------
  \coordinate (bg) at ([yshift=-4.8mm]$(g1.south)!0.5!(g2.south)$);
  \node[human, pfull] (ap) at (0,0 |- bg)
    {\textbf{Operator \mode{Apply}}\ \ an explicit human action over the diff and
     the evaluation --- never automatic};
  \draw[flow] (g1.south) -- (g1.south |- ap.north);
  \draw[weak, densely dashed, draw=cbGovern] (g2.south) -- (g2.south |- ap.north);

  \node[extern, pfull] (pr) at ([yshift=-3.4mm]ap.south west)
    {\textbf{\mode{Promote}}\ \ audited alias rotation; the outgoing target is
     recorded first, so rollback restores (\algref{alg:promote})};
  \draw[flow] (ap.south) -- (pr.north);

\end{tikzpicture}
\end{cbscaled}
\caption{Inside one refinement job --- the resolution at which the gate semantics
become visible. Evidence assembly folds trace evidence and versioned examples
into one corpus; diagnosis is persisted as an artifact rather than kept in a
prompt; a policy selects among \statOptimizers{} implemented optimizer paths; and
the resulting candidate is scored per dimension against the incumbent. The two
gates are drawn differently because they are different mechanisms, and the paper
depends on the distinction. The heavy-bordered \advgate{} is \emph{enforced}: it
decides whether the job may advance to \code{candidate\_ready}, and a failure
stops the job. The dashed \relverdict{} is persisted per version as release
evidence and never blocks alias rotation. Between them sits a human: \mode{Apply}
is an explicit operator action over the diff and the scores, not a threshold that
fires on its own.}
\label{fig:pipeline}
\end{figure}


Calibration is the machinery that proposes a measurably better version, and
\figref{fig:pipeline} shows one job's internals. \statOptimizersC{} optimizer
paths are implemented: two exposed on the prompt form, two selectable by policy, and one
requiring explicit configuration. Non-prompt families use their own idioms ---
manifest replay for workflows, revision-fenced deterministic suites for tools,
retrieval-quality calibration for knowledge bases.

Two properties of this component are architectural rather than algorithmic, and
they are the ones that matter for the paper's argument. First, the diagnosis is
persisted as a durable artifact rather than held in a prompt, so the reasoning
behind a candidate outlives the job that produced it. Second, optimizer selection is
a \emph{policy} decision recorded on the job, so a later reader can tell which
strategy produced a given candidate --- without which a comparison across candidates
is uninterpretable.

\subsection{Gate semantics: the distinction the system turns on}

%% alg-gate.tex -- Algorithm: the two gates, side by side.
%%
%% Written as one figure precisely because the two mechanisms are constantly
%% conflated. Reading the two procedures together makes the asymmetry structural
%% rather than rhetorical: one returns a decision that stops the job, the other
%% returns a record that is filed.

\begin{algorithm}[tbp]
\caption{The two gates, side by side. \code{advance} \emph{stops the job}: its
return value is control flow, and there is no override path because withholding a
candidate that failed its own evaluation is never what an operator needs to
override. Read ``unbypassable'' as scoped to this state machine: a direct authoring
route can move a prompt alias without entering it (\secref{sec:limitations}). \code{verdict} \emph{files a record}: its return value is evidence,
consumed by a human at \mode{Apply}. Placing the unbypassable check on advancement
rather than on release is design decision D3 (\tabref{tab:decisions}).}
\label{alg:gate}
\psig{advance}\ \ \cmtin{ENFORCED --- the return value is control flow}
\pin{$v$}{a candidate; $b$, the incumbent baseline; $C$, the assembled corpus}
\pout{\plit{candidate\_ready} \kw{or} \plit{stopped}}{}
\begin{pseudo}
\pline{$S \gets \emptyset$}
\pline{\kw{for each} \pid{scorer} \kw{in} scorecard($v$) \pdo}
\pline{\I \kw{for each} \pid{x} \kw{in} $C$ \pdo}
\pline{\II $S[\pid{scorer}] \gets S[\pid{scorer}] \cup
        \{$\pid{scorer}$(v, x)\}$}
\pline{}
\pline{\kw{for each} \pid{d} \kw{in} dimensions($S$) \pdo}
\pline{\I \kw{if} mean($S[d]$) $<$ mean(baseline($b$, $d$)) $- \epsilon_d$
        \pthen}
\pline{\II persist($v$, $S$, \plit{stopped}, reason $= d$)}
\pline{\II \kw{return} \plit{stopped}
        \cmt{no human ever sees $v$}}
\pline{}
\pline{persist($v$, $S$, \plit{candidate\_ready})}
\pline{\kw{return} \plit{candidate\_ready}}
\end{pseudo}

\vspace{0.8ex}
\psig{verdict}\ \ \cmtin{ADVISORY --- the return value is evidence}
\pin{$v$}{a version, already released or releasable}
\pout{a persisted per-version record}{never a control-flow decision}
\begin{pseudo}
\pline{$R \gets \langle$ scores($v$), judges($v$), corpus id, timestamp
       $\rangle$}
\pline{persist\_version\_verdict($v$, $R$)}
\pline{\cmtin{deliberately absent: any branch on $R$ that could block}}
\pline{\cmtin{rotate(). A stale verdict must never gate an urgent fix.}}
\pline{\kw{return} $R$ \cmt{surfaced in review}}
\end{pseudo}
\end{algorithm}


\sys has two mechanisms that a reader may be tempted to call ``the gate,'' and
conflating them would invalidate several claims in this paper.

\para{The \advgate{} is enforced.} It decides whether a refinement job may advance
to a candidate-ready state. A failing gate stops the job. Nothing downstream sees
the candidate; in particular, no operator is asked to review it. Within the
refinement state machine this gate cannot be bypassed --- but that scope is the
whole claim. A direct operator release uses the durable release service but may
carry an advisory verdict and attributed override without entering the machine.

\para{The \relverdict{} is advisory.} It is persisted per version as release
evidence and surfaced in review. It does \emph{not} block prompt alias rotation.

\para{The \depgate{} is family-specific.} On asset paths that wire it --- workflow
deployments most notably --- a deploy-gate policy enforces release under an
optimistic alias check. It is not a cross-family policy.

The reason for placing the enforced gate on advancement rather than on release is
the subject of D3 in \tabref{tab:decisions}, and is worth restating because it is
counterintuitive. A gate that blocks release must be overridable, because verdicts
go stale, corpora drift, and an urgent fix must be shippable at 3am. Once a gate is
overridable, the override becomes the normal path, and the gate stops carrying
information. Placing the unbypassable check earlier --- where its only power is to
withhold a candidate that failed its own evaluation --- keeps it meaningful, because
withholding a bad candidate is never the thing an operator needs to override.

\algref{alg:gate} sets the two procedures side by side, which makes the asymmetry
structural rather than rhetorical. \code{advance} returns control flow: its
\code{stopped} branch ends the job and no human is ever shown the candidate.
\code{verdict} returns a record: it persists evidence and, by construction,
contains no branch that could block a rotation. The absence on lines~3--4 of that
second procedure is the design, not an omission.

\subsection{Release control}

%% alg-promote.tex -- Intent-first prompt alias release and reconciliation.

\begin{algorithm}[tbp]
\caption{Intent-first release for a prompt alias. The exact outgoing and incoming
versions, actor, scopes, and evidence are committed before the provider call. A
crash after the call therefore leaves an observable \plit{applying} operation,
not an unrecorded production change. Alias assignment to the same version is
idempotent, so the operation identifier is also the retry key. This algorithm is
implemented for prompt alias promotion and rollback; other families retain the
weaker guarantees in \tabref{tab:families}.}
\label{alg:promote}
\psig{release}
\pin{$a$}{prompt asset; $t$, alias; $v$, immutable version; $p$, principal}
\pin{$o$}{caller-supplied or generated release operation identifier}
\pout{applied}{or a durable reconciliation obligation}
\begin{pseudo}
\pline{\kw{if} release[$o$] exists \pthen{} \kw{return} retry(release[$o$])}
\pline{\kw{if} \plit{operator} $\notin$ scopes($p$) \kw{and}
       \plit{admin} $\notin$ scopes($p$) \pthen{} \kw{reject}}
\pline{$u \gets$ resolve($a,t$) \cmt{exact outgoing version}}
\pline{\kw{begin transaction}}
\pline{\I release[$o$] $\gets \langle$\plit{prepared}, $a,t,u,v,p,$ scopes,
       evidence$\rangle$}
\pline{\I audit(\plit{prepare}, $o,a,t,u,v,p$)}
\pline{\kw{commit} \cmtin{the recovery record now survives a crash}}
\pline{}
\pline{\kw{begin transaction}; release[$o$].status $\gets$ \plit{applying};
       \kw{commit}}
\pline{\kw{try} rotate($a,t \to v$) \cmt{idempotent provider effect}}
\pline{\kw{catch} $e$}
\pline{\I release[$o$] $\gets \langle$\plit{reconcile\_required}, $e\rangle$;
       \kw{commit}; \kw{return}}
\pline{\kw{begin transaction}}
\pline{\I release[$o$] $\gets \langle$\plit{applied}, provider result,
       \pid{now}()$\rangle$}
\pline{\I audit(\plit{promote}, $o,a,t,u,v,p$)}
\pline{\kw{commit}}
\end{pseudo}

\vspace{0.5ex}
\psig{reconcile}
\pin{$o$}{an \plit{applying} or \plit{reconcile\_required} operation}
\begin{pseudo}
\pline{$x \gets$ resolve(release[$o$].asset, release[$o$].target)}
\pline{\kw{if} $x =$ release[$o$].version\_after \pthen{}
       settle($o$, \plit{applied})}
\pline{\kw{else if} $x =$ release[$o$].version\_before \pthen{}
       settle($o$, \plit{failed})}
\pline{\kw{else} settle($o$, \plit{reconcile\_required})
       \cmt{unexpected third state}}
\end{pseudo}
\end{algorithm}


Release is where the audit obligations concentrate. Prompt creation and version
authoring now register immutable drafts and cannot rotate an alias. Every prompt
alias promotion and rollback enters the state machine in \algref{alg:promote}:
CALIBER commits the exact outgoing and incoming versions, actor, scopes, evidence,
and an operation identifier before calling \mlflow{}. It then marks the operation
\code{applying}, performs an idempotent alias assignment, and settles it as
\code{applied}; an exception becomes \code{reconcile\_required}. The release API
accepts the operation identifier on retries, and an operator-visible reconciler
settles incomplete rows from the alias actually observed in \mlflow{}.

This closes the unrecorded-effect window for prompt aliases, not for every family.
Workflow deployments maintain a checkpoint stack that rollback pops; knowledge-base
activation derives the prior active build from history; skill rollback restores a
prior snapshot as a \emph{new} current version, which is a different guarantee and
is listed as such.

\para{The audit guarantee, stated exactly.} For a mutation whose effects are
entirely database-resident and which calls the audit helper on the caller's
session, the mutation and its audit row share one SQL transaction and cannot
diverge. That is the whole claim. It does not extend to external \mlflow{} or
provider effects, which are dual-write boundaries; and it is a property of a call
site, not of a module, so importing the helper proves nothing about coverage. We
state it this narrowly on purpose: an unqualified ``everything is audited'' claim
would be unfalsifiable, and the qualified one is checkable.

\subsection{The trust boundary}

%% tab-threat.tex -- Table: the threat model, stated as a table rather than prose.
%%
%% A reviewer asked for this specifically, and the request was well founded: the
%% security discussion was spread across three sections, which let a reader
%% assemble a stronger model than the one that holds. A table forces every row to
%% name its residual risk and its evidence, and the evidence column is the point
%% -- most rows are argued by construction and none is measured, which is far more
%% obvious here than it was in prose.
%%
%% "Containment" and "isolation" are distinct terms throughout this paper. The
%% distinction is load-bearing and this table is where it is most visible.

\begin{table}[tbp]
\centering
\caption{Threat model. Each row names an adversary, what must be trusted for that
row's boundary to mean anything, the asset it protects, the boundary actually
enforced, what remains exposed, and where the claim can be checked. The evidence
column is the honest one: every row is argued \emph{by construction}, and
\textbf{no row is supported by measurement} --- \secref{sec:eval} has no adversarial
experiment, and a reader should not read this table as a security evaluation. Rows
1--3 are where \sys is weakest, and the containment-versus-isolation distinction
they turn on is used consistently throughout (\figref{fig:trust}).}
\label{tab:threat}
\small
\setlength{\tabcolsep}{4pt}
\renewcommand{\arraystretch}{1.10}
\begin{cbwide}
\begin{tabular}{@{}L{0.30cm}L{2.28cm}L{2.02cm}L{2.16cm}L{3.06cm}L{3.26cm}L{2.18cm}@{}}
\toprule
\theadrow
\thead{} & \thead{Adversary} & \thead{Must be trusted} & \thead{Protected asset} &
\thead{Boundary enforced} & \thead{Residual risk} & \thead{Evidence} \\
\midrule

1 & Author of a local tool &
the tool author &
host filesystem, network, secrets &
separate interpreter, empty environment, private working directory, resource
limits &
\textbf{containment, not isolation}: ambient host filesystem and network authority
remain &
by construction (\secref{sec:components}) \\

\zebra
2 & Local \code{stdio} MCP server &
the server binary &
the same &
command and host allowlists, sanitized environment, private cwd, process timeout &
same as row~1; the shipped bubblewrap profile is containment~\cite{bubblewrap} &
by construction \\

3 & Remote MCP endpoint &
\chipnone{} &
egress, tool inventory &
fail-closed egress policy; production preflight demands an attested sidecar or
non-loopback HTTPS &
policy covers \emph{governed} transports only, not arbitrary user code &
by construction \\

\zebra
4 & Retrieved content (indirect prompt injection) &
nothing &
tool authority reachable from a run &
capability risk tiers; \tier{gated} capabilities withheld from the synchronous
projection &
protection is only as strong as the classification, and two \tier{mutate} tools
approve and publish &
\chipunmeasured{}\newline\cite{greshake2023injection} \\

5 & Caller of a published service &
\chipnone{} &
enqueue capacity, workflow authority &
Bearer validation before the body is read, body cap, quota, idempotency; policy
revalidated under the enqueue lock &
a per-request application bound, \textbf{not} connection-level flood protection &
by construction \\

\zebra
6 & Account client escalating scope &
the session store &
live targets &
one scope lattice; \operator{} and \approver{} sibling, \admin{} implies all
\statScopes{} &
no separation of duties: one operator may author and apply
(\secref{sec:limitations}) &
by inspection \\

7 & Holder of \mlflow{} registry credentials &
the registry's own access control &
the prompt live target &
\chipnone{} --- outside \sys's control &
an alias moved directly leaves no intent record and no evidence; \sys's guarantees
are about \emph{its} call sites &
by construction (\secref{sec:limitations}) \\

\bottomrule
\end{tabular}
\end{cbwide}
\tabnote{``By construction'' means the boundary follows from the code path and is
checkable by reading it; ``by inspection'' means it was verified at the call sites
but is not enforced by a type or a test. Neither is a measurement. Closing row~1
and row~2 requires an OS-enforced backend~\cite{agache2020firecracker,bubblewrap},
which \secref{sec:future} records as the platform's sharpest outstanding security
gap.}
\end{table}


\figref{fig:trust} gives the admission and authority path, and
\tabref{tab:threat} states the threat model row by row --- adversary, what must be
trusted, the asset, the boundary actually enforced, the residual risk, and where the
claim can be checked. We put it in a table rather than in prose because prose let a
reader assemble a stronger model than the one that holds: the evidence column makes
it plain that every row is argued by construction and none is measured. Two entry classes have
genuinely different admission rules: account clients present a session (or, in
opt-in trusted-header proxy mode, a header) while published-service callers present
either a public policy or a Bearer token, with a body cap, a quota, and
idempotency handling. Both converge on one scope lattice in which \operator{} and
\approver{} are sibling scopes each implying \viewer{}, and \admin{} implies all
\statScopesW{}.

Two properties need stating carefully, and both are weaker than a casual reading
would suggest.

First, \aria{} capabilities are risk-tiered. Registry capabilities declared
\tier{gated} are withheld from the synchronous tool projection and require a plan
interaction. The hand-written draft approval and publication tools are also
classified \tier{gated}; even a permissive build-mode session neither advertises nor
dispatches them. The manual path uses the draft lifecycle UI/API. Two additional
draft policies cross a separate service boundary: \code{agent\_review} presents an
immutable candidate hash, validation report, and test report to an isolated
no-tools reviewer agent whose service identity must differ from the author and hold
\approver{}; \code{full\_autonomy} continues an approved decision through a third,
\operator{}-scoped release identity. The decision records policy/model version,
rationale, confidence, evidence identifiers, and expiry; invalid output, missing
scope, shared identity, low confidence, expiry, or candidate drift fails closed.
Goal-plan \tier{gated} interactions remain human decisions. These gated paths are
not yet one unified registry mechanism. The protection remains only as
strong as future classifications, which is why the registry projection and its
contract tests matter.

Second, local stdio MCP controls are \emph{containment}, not a sandbox. Command and
host allowlists, a sanitized environment, a private working directory, and process
timeouts reduce ambient authority. Only an attested managed sidecar or a
non-loopback HTTPS endpoint counts as an execution boundary for production
preflight; the shipped bubblewrap profile is containment~\cite{bubblewrap}.
Similarly, fail-closed egress applies to governed HTTP and MCP transports and not
to arbitrary user code: the shipped local tool runner has a separate interpreter,
an empty environment, a private working directory, and resource limits, but retains
ambient host filesystem and network authority. Untrusted authors require an
operator-supplied OS-enforced backend~\cite{agache2020firecracker}. Throughout this
paper, ``containment'' and ``isolation'' are used as distinct terms, and the
distinction is load-bearing rather than stylistic.

\subsection{Published services and the copilot}

A workflow may be published as an HTTP service under two-stage admission: a
protected invoke validates its Bearer token before the body is consumed, and both
public and protected invokes then count raw request chunks against a configured
maximum before parsing, with policy and token revalidated under the enqueue lock.
This is a per-request application bound and explicitly not
connection-level flood protection --- a distinction worth preserving, because
conflating the two is how a body cap gets mistaken for a denial-of-service defence.

\aria{}, the embedded copilot, is a permissioned agentic tool loop with durable
goal-plans and asynchronous resume: plans parked on long-running jobs are picked up
by the plan worker when those jobs settle. It is a client of the same capability
and authority model as every other surface, which is the property that keeps it
from becoming a privilege-escalation path.
